% iononshpere/main.tex
% English arXiv-style manuscript.

\documentclass[11pt,a4paper]{article}
\usepackage[a4paper,margin=1in]{geometry}
\usepackage{amsmath,amssymb,bm}
\usepackage{booktabs}
\usepackage{array}
\usepackage{hyperref}
\usepackage{siunitx}

\newcolumntype{L}[1]{>{\raggedright\arraybackslash}p{#1}}

\setlength{\emergencystretch}{5em}

\hypersetup{
  colorlinks=true,
  linkcolor=blue,
  urlcolor=blue,
  citecolor=blue,
  pdftitle={Dual Asymmetric Constraints and Stratified Cascade in Planetary Ionospheres},
  pdfauthor={Ting Peng},
  pdfkeywords={ionosphere; plasma; hypothesis; non-symmetric constraints; gravity; magnetic field; light--heavy ion collisions; cascade; corona}
}

\title{Dual Asymmetric Constraints and Stratified Cascade in Planetary Ionospheres:\\
A Working Hypothesis for Elevated Ionospheric and Coronal Temperatures}

\author{Ting Peng\\
Key Laboratory for Special Area Highway Engineering of Ministry of Education,\\
Chang'an University, Xi'an 710064, China\\
\url{t.peng@ieee.org}\\
ORCID: \href{https://orcid.org/0009-0001-9059-2278}{0009-0001-9059-2278}
}

\date{\today}

\begin{document}

\maketitle

\begin{abstract}\raggedright\sloppy
A working hypothesis is advanced whereby several ionospheric and coronal phenomena are treated within a single kinetic--geometric framework. The account addresses elevated effective temperatures in ionospheres and tenuous coronae relative to the underlying dense surface or photosphere; the frequent increase of effective temperature with altitude; preferential accumulation of the lightest abundant ionic species (notably H$^+$) in the outermost layers when present; the dependence of the highest attainable suprathermal-tail temperatures on the vertical distribution of ion masses; and systematic differences between planetary environments---for instance, comparatively weak Venus ionospheric structure under heavy-ion dominance, limited light-ion abundance for multi-step energization, and weak global magnetic confinement, as opposed to comparatively moderate topside thermal enhancement in hydrogen--helium-dominated ionospheres such as Jupiter's, where the availability of heavy collision partners for repeated mass-asymmetric transfer may be restricted. The structural elements common to these applications are gravitational speed-space selection, multi-stage elastic collisions between light and heavy ions that redistribute kinetic energy upward along the column, and magnetic geometry where it is present, which favors maintenance of stratified structure and sustained cross-scale coupling relative to rapid scattering-driven dilution.

The thermal component of the proposal is presented as falsifiable rather than as a closure validated against full radiative--collisional ionospheric models. In parallel, an electrodynamic thread is developed: asymmetric constraints bias steady occupation away from Boltzmann form and can support directed transport without macroscopic thermal or chemical-potential gradients; gravity acts as a speed-space filter while the magnetic field rectifies seed charge separation into a current, and a stratified D/E/F representation implements layer-resolved electrostatic cascade amplification. An order-of-magnitude closure employs an illustrative terminal electric-field amplitude $\sim 10^2\ \mathrm{V\,m^{-1}}$ as a dimensional anchor for the algebraic structure; this magnitude is not identified with typical large-scale ambient ionospheric DC fields ($\mathrm{mV\,m^{-1}}$--$\mathrm{V\,km^{-1}}$). The idealized formulation omits explicit solar heating and photoionization; quantitative comparison with observations requires restoration of those sources and kinetic calibration of the thermal claims. In summary, this paper advances a gravity-driven, multi-component cascade energization and heating hypothesis that provides a unified physical picture for anomalously high temperatures in both the corona and the ionosphere; its validity remains to be tested by future observations.
\end{abstract}

\noindent\textbf{Keywords:} working hypothesis; non-equilibrium steady state; non-symmetric constraints; ionospheric electrodynamics; gravity--magnetic rectification; light--heavy ion collisions; stratified cascade; electrostatic closure; corona.

\section{Introduction}
\label{sec:intro}
\subsection*{Working hypothesis and scope}
References to a ``mechanism'' in what follows denote a \emph{testable working hypothesis} whose purpose is to relate several observations that are often treated independently: the increase of effective ionospheric temperature with altitude, the high effective temperatures of tenuous coronae above cooler photospheres, species-dependent vertical stratification (with the lightest ions preferentially found in the outermost layers when abundant), and systematic contrasts between environments differing in ion mass spectra and magnetic field strength. The central proposition is that gravitational effects sort faster particles to greater altitudes, that multi-stage elastic collisions between light and heavy ions mediate repeated upward redistribution of kinetic energy along the column, and that magnetic fields, where sufficiently strong, inhibit rapid loss of coherent layered structure through scattering and cross-field transport. The electrostatic cascade developed in subsequent sections is offered as a complementary pathway---charge separation followed by amplification---that may coexist with the thermal description; both threads invoke the same pair of constraints (gravity and magnetism) but emphasize distinct observables (electric fields and currents as opposed to suprathermal tails and composition).

Classical atmospheric and ionospheric theories are typically built on the assumption of \emph{external energy input}: solar UV/X-ray ionization provides the charge carriers, solar heating provides the thermal energy, and temperature gradients drive large-scale circulation. Yet a set of observational features remains difficult to reconcile within a single traditional picture:
\begin{enumerate}
  \item ionospheric temperatures increase with altitude, contrary to the expected thermal stratification;
  \item even during local night when solar radiation is absent, the ionosphere can remain active, maintaining stable electron density and elevated temperatures;
  \item electron temperatures can significantly exceed ion temperatures, accompanied by persistent charge separation;
  \item large-scale current systems (e.g., equatorial and auroral current systems) can persist without an obvious, continuous external electrical power source;
  \item in plasmas, electron and ion motions can exhibit strong directional asymmetry, with velocity--altitude decoupling.
\end{enumerate}
These features, including electron--ion temperature decoupling and structural stratification in conductivity and current systems, are standard topics summarized in ionospheric textbooks and reviews \cite{Kelley2009EarthIonosphere,SchunkNagy2009IonospherePhysics}.

Recent statistical physics studies of constraints \cite{peng2026entropydirectionmirrorstateparadox} provide a new principle: \emph{non-symmetric constraints can alter the occupation distribution so that it deviates from the Boltzmann form; directed transport can arise even without temperature differences or chemical-potential gradients}.

The principle is applied here to planetary ionized envelopes. In the absence of macroscopic temperature gradients, the combination of gravity and magnetic field constitutes a natural dual asymmetry: the constraints rectify thermal motion into a weak initial directed charge flow and seed electric field. The multi-layer ionospheric structure then permits \emph{cascade amplification} of that seed signature. Relating the resulting electrostatic amplitude to specific observables requires a precise definition of the field variable $E$ and, in general, rescaling to reported ionospheric magnitudes (Sec.~\ref{sec:efield-scales}) \cite{Kelley2009EarthIonosphere,SchunkNagy2009IonospherePhysics}.

The phrase ``no external driving'' is used in the sense of absent continuous non-thermal energy injection; the static constraints function as filters and rectifiers that reorganize thermal energy into electrostatic structure.

\subsection*{Contributions}
\begin{enumerate}
  \item A working hypothesis is formulated for elevated temperatures in ionospheres and coronae, combining gravitational speed-space filtering, multi-step light--heavy ion collisional transport of kinetic energy upward along the column, and magnetic stabilization of stratification where the field is strong.
  \item A dual-asymmetry ionospheric formulation is presented in which gravity acts in speed space and magnetic geometry rectifies momentum transport.
  \item A stratified cascade description (D/E/F1/F2) relates seed charge separation to macroscopic electrostatic signatures through inter-layer coupling.
  \item An explicit electrostatic closure is constructed with effective parameters $(\eta,\xi,\zeta)$ and multiplicative layer gains.
  \item Scope conditions are stated: closure amplitudes are illustrative; comparison with observations requires external source terms and transport physics; cross-object trends (Venus, Jupiter, hydrogen-dominated topside) are advanced as qualitative predictions subject to test against composition and magnetic data.
\end{enumerate}

\section{Conceptual tensions in canonical descriptions}
Canonical ionospheric and atmospheric descriptions exhibit a number of internal tensions:
\begin{enumerate}
  \item \textbf{Altitude temperature anomaly:} standard thermodynamic stratification suggests temperature should decrease with altitude, yet ionospheric temperature often rises.
  \item \textbf{Charge separation without an obvious external driver:} photoionization vanishes in the night, while ionospheric signatures do not disappear.
  \item \textbf{Strong electron--ion asymmetry:} due to mass disparity, stable velocity ratios can be far from what naive thermal equilibrium would suggest.
  \item \textbf{Directed transport without strong temperature gradients:} persistent current loops can exist without a dominant temperature-gradient drive.
  \item \textbf{Field-amplitude mismatch:} reconciliation of microscopic seed charge separation with macroscopic electrostatic signatures appears to require an intermediate amplification mechanism (Sec.~\ref{sec:efield-scales} specifies the treatment of amplitudes adopted here).
 \end{enumerate}
Collectively, these considerations motivate an intrinsic, constraint-driven ordering process that incorporates \emph{micro-to-macro amplification}, schematically represented here by multi-layer cascade effects.

\section{Theoretical framework: directed transport and cascade amplification under dual asymmetry}

\subsection{Statistical-physics foundation of non-symmetric constraints}
The key conclusions of \cite{peng2026entropydirectionmirrorstateparadox} can be summarized as:
\begin{enumerate}
  \item constraint asymmetry breaks detailed balance and drives the steady state away from the Boltzmann distribution;
  \item directed particle flow can be generated directly by constraint structure without requiring temperature difference or chemical-potential gradients;
  \item spatial anisotropy of constraints itself acts as a drive for directed transport;
  \item weak directed transport can be hierarchically amplified through structural coupling, producing macroscopic observable effects.
\end{enumerate}
In the ionospheric implementation adopted here, cascade amplification is not introduced as an external energy source; it follows from representing the ionosphere as a hierarchy of coupled response channels. Translation of the abstract statistical-mechanical constraints of \cite{peng2026entropydirectionmirrorstateparadox} into fluid or kinetic plasma models requires further specifications (collision operators, boundary conditions, magnetization). The ionospheric block is therefore presented as a physically motivated realization rather than as a formal corollary of the abstract framework.

\subsection{Gravity as a speed-space asymmetry (seed charge separation)}
Let $\phi(r)$ denote gravitational potential. In the vertical direction, gravity introduces an energy asymmetry:
\begin{itemize}
  \item upward motion requires overcoming the gravitational potential increment, effectively filtering the thermal kinetic energy that can reach higher altitudes;
  \item downward motion releases mechanical energy as potential decreases, without requiring any external energy input.
\end{itemize}

\subsubsection{Model object and assumptions}
The analysis addresses thermal carriers in the ionosphere (electrons and dominant ions) within a reduced gravity--thermal model specified by the following assumptions:
\begin{enumerate}
  \item altitude domain: from $\sim \SI{60}{km}$ (bottom D layer) to $\gtrsim \SI{1000}{km}$ (top F2 region), with gravity approximated as a radially uniform field $g\approx \SI{9.8}{m/s^2}$ (high-altitude corrections can be implemented);
  \item particle species: electrons and representative ions (e.g., protons and oxygen ions), neglecting neutral-collision energy losses at the dominant scale (with appropriate corrections at lower altitudes);
  \item dynamics: carriers feel gravity and thermal motion only; initial thermal velocities follow an isotropic Maxwellian distribution;
  \item filtering criterion: only carriers whose thermal kinetic energy exceeds the gravitational potential increment can reach the upper regions and become effective seeds for charge separation and cascade amplification.
\end{enumerate}

\subsubsection{A quantitative estimate of high-altitude speed}
By energy conservation, for a carrier moving from $h_1$ to $h_2$:
\begin{equation}
\frac{1}{2}mv_1^2+mgh_1=\frac{1}{2}mv_2^2+mgh_2,
\end{equation}
so that
\begin{equation}
v_2=\sqrt{v_1^2-2g\Delta h},
\qquad \Delta h=h_2-h_1.
\end{equation}
Using representative ionospheric altitudes $h_1=\SI{60}{km}$ and $h_2=\SI{300}{km}$, $\Delta h=\SI{240}{km}=\SI{2.4e5}{m}$. For electrons we use a representative thermal scale $v_{1e}\approx \sqrt{2kT_e/m_e}$ with an elevated F-region electron temperature $T_e\sim 10^4~\mathrm{K}$, which yields $v_{1e}\sim 5\times 10^5~\mathrm{m/s}$ (consistent with $T_e\gg T_i$ in ionospheric physics \cite{Kelley2009EarthIonosphere,SchunkNagy2009IonospherePhysics}); we adopt $v_{1e}\approx 5.5\times 10^5~\mathrm{m/s}$ as a concrete estimate. For ions we take a proton scale $v_{1i}\approx 1.3\times 10^4~\mathrm{m/s}$ (near-thermal at $\sim 300~\mathrm{K}$). Then $v_{2e}\approx v_{1e}$ (negligible reduction) and $v_{2i}\approx 1.28\times 10^4~\mathrm{m/s}$ (small reduction).

\subsubsection{Persistence of elevated kinetic energy after filtering}
The following considerations apply:
\begin{enumerate}
  \item the thermal speed scale is far larger than the speed loss $\sqrt{2g\Delta h}$ required to climb $\Delta h$;
  \item filtering is \emph{selective} rather than directly \emph{accelerating}: only fast carriers survive the threshold condition, so the surviving population remains high-energy;
  \item at sufficiently high altitude, collision losses are weak enough that the remaining kinetic-energy content is not strongly dissipated.
\end{enumerate}
Indirect processes (e.g., collision-mediated redistribution and electrostatic field-driven directed acceleration during cascade) can further enhance directional motion while preserving energy bookkeeping.

\subsubsection{Gravitational filtering as selection rather than acceleration}
Within the energy-conservation formulation above, gravitational filtering constitutes a \emph{selection} mechanism, not an \emph{acceleration} mechanism in the sense of net mechanical-energy gain from the gravitational field alone. The mapping from low-altitude to high-altitude kinetic energy amounts to exchange between kinetic and gravitational potential energy under a conservative force and does not increase total mechanical energy above its initial thermal value.

Two \emph{indirect} channels may nevertheless produce macroscopic signatures resembling acceleration:
\begin{enumerate}
  \item \emph{Collision-mediated redistribution (predominantly for electrons).} Elastic and inelastic encounters between fast electrons and residual neutrals can redistribute kinetic energy among species, yielding transient increases in selected particles' speeds. This process reallocates energy within the thermal budget rather than injecting power from an external source.
  \item \emph{Electrostatic acceleration coupled to cascade amplification.} Gravitational filtering seeds weak charge separation and hence a seed electric field. As the cascade proceeds, the macroscopic field may grow in successive layers. Carrier acceleration along the field is then driven by the amplified electrostatic structure and is contingent on that electrodynamic evolution, not on gravity acting as a direct heating source.
\end{enumerate}
Thus gravitational filtering does not, by itself, raise carrier speeds above the initial thermal distribution; indirect signatures of enhanced directed motion may nonetheless arise from the conjunction of selective filtering and subsequent field amplification, with electrons typically exhibiting stronger effects than ions.

\subsubsection{Cyclic exchange of kinetic and potential energy}
A closed sequence consistent with energy conservation may be formulated that sustains elevated high-altitude kinetic-energy content in idealized circumstances without continuous external driving:
\begin{enumerate}
  \item \emph{Filtering and collisional repopulation of the fast population.} Gravitational filtering transports a subpopulation of high-speed carriers to greater altitudes. In intermediate collisional regimes (for example the E region and adjacent altitudes), elastic collisions redistribute kinetic energy from initially fast to initially slower carriers, repopulating the high-speed tail. When multiple ion species are present, successive encounters between light and heavy ions constitute a \emph{multi-stage ladder}: heavy species may serve as momentum reservoirs such that light ions, in favorable collision geometries, gain speed across repeated interactions and become more likely to satisfy upward filtering criteria. This sequence supplies the kinetic basis for the thermal working hypothesis stated in Sec.~\ref{sec:intro}.
  \item \emph{Energy loss and gravitational settling.} Carriers whose kinetic energy falls below the filtering threshold are not gravitationally supported at high altitude and settle downward.
  \item \emph{Recovery of kinetic energy on descent.} During descent, gravitational potential energy converts to kinetic energy; at lower altitudes, speeds may again exceed the filtering threshold, completing the cycle.
\end{enumerate}
The cycle involves only interchange between kinetic and gravitational potential energy; its structural admissibility is therefore not precluded by omission of radiative terms from the idealized core. When collisional dissipation is weak on the relevant timescales, such a loop may sustain a supply of high-energy carriers compatible with cascade amplification and with persistent departures from naive hydrostatic temperature profiles, including day--night persistence in schematic models.

\subsection{Magnetic field as directional asymmetry (seed current formation)}
The Lorentz force,
\begin{equation}
\vec F = q\vec v\times \vec B,
\end{equation}
introduces directional asymmetry:
\begin{itemize}
  \item along $\vec B$, the Lorentz force vanishes and motion along the field line is unconstrained by magnetism at leading order;
  \item perpendicular to $\vec B$, the Lorentz force enforces gyromotion and inhibits free cross-field displacement over large distances.
\end{itemize}
Because the force reverses sign between species, electrons and ions acquire distinct drift and transport patterns. In the absence of macroscopic temperature gradients, magnetic asymmetry converts seed charge separation into a directed current, which subsequently drives cascade amplification in the layered model.

\subsection{Cascade amplification in the stratified ionosphere}
The ionosphere is effectively partitioned by altitude into D, E, F1, and F2 regions. Each region features distinct conductivity and different collision regimes. The cascade amplification is then a hierarchy of coupled steps:
\begin{enumerate}
  \item D layer seeds weak charge separation and initiates a seed current;
  \item E layer increases effective current transmission and yields enhanced electric-field superposition;
  \item F1 and F2 layers provide the terminal amplification through strong conductivity and weak collision losses, allowing charge accumulation and macroscopic field formation.
\end{enumerate}

\subsubsection{A detailed three-step cascade mechanism}
The core cascade is a stepwise, bottom-to-top, positive-feedback process. For the order-of-magnitude \emph{closed} estimate, it can be summarized in three layer-to-layer steps (a qualitative mechanism consistent with the electrostatic parameterization):

\paragraph{Step 1: D layer (seed excitation layer) --- weak charge separation and current initiation}
The D region is the lowest ionospheric layer, characterized by neutral dominance, frequent collisions, and comparatively weak constraint strength. It nevertheless serves as the initiation layer in the schematic cascade:
\begin{enumerate}
  \item \emph{Gravitational filtering:} a small fraction of energetic electrons reaches D-region altitudes; collisions with neutrals supply a weak ionization source (in the idealized limit without solar radiation, fast-particle collisions still provide a minimal carrier population at the seed stage).
  \item \emph{Charge separation:} anisotropic electron motion relative to ion response yields an initial weak electric field $E_1\approx E_0\sim 8\times10^{-5}\ \mathrm{V/m}$.
  \item \emph{Seed current:} magnetic rectification of the separation yields a seed current $J_1\approx J_0\sim 8\times10^{-10}\ \mathrm{A/m^2}$.
  \item \emph{Upward coupling:} the weak current and field constitute the seed signal passed to higher layers; low D-region conductivity limits local accumulation and favors coupling into the E region.
\end{enumerate}

\paragraph{Step 2: E layer --- current enhancement and field superposition}
E has higher conductivity than D, lower collision frequency, and a stronger magnetic constraint, enabling the first significant amplification:
\begin{enumerate}
  \item \emph{Charge injection:} the directed electron flow from D injects carriers into E, and additional collisions increase the number of charged particles.
  \item \emph{Current amplification:} effective constraint duration increases and ionization generation improves, enhancing the directed current density by about an order of magnitude, i.e. an effective factor $\mathcal{A}_{\mathrm{E}}\approx 10$ in the electrostatic closed sense.
  \item \emph{Electric-field superposition:} enhanced charge separation leads to a superposition of fields. At the level of the target order-of-magnitude estimate, $E_2+E_1\sim 8\times10^{-4}\ \mathrm{V/m}$, i.e. roughly another factor $\sim 10$ relative to $E_1$.
  \item \emph{Positive feedback:} the enhanced field drives additional directed motion, increasing the current and sustaining the feedback loop.
\end{enumerate}

\paragraph{Step 3: F1 + F2 layers --- charge accumulation and macroscopic field formation}
F1 and F2 are the terminal regions, where carrier density and conductivity are largest and collisions are least frequent (F2 is nearly collisionless):
\begin{enumerate}
  \item \emph{Sustained current injection:} the amplified transport from E injects into F1, and in the electrostatic closed model this coupling corresponds to a large effective field gain $\mathcal{A}_{\mathrm{F1}}\approx 1000$.
  \item \emph{Charge accumulation:} in F2, electrons accumulate with minimal collisional dissipation; meanwhile positive charge can pile up near the interface with F1 due to inertia/response delay of ions.
  \item \emph{Final field build-up:} superposition of layer fields, combined with very high F2 conductivity, strengthens charge separation and yields the macroscopic field $E_{\mathrm{total}}\sim 80\ \mathrm{V/m}$.
  \item \emph{Stability:} with minimal collisional energy dissipation, the charge separation and field can be maintained; in the ideal limit, thermal replenishment continues to sustain the amplification chain.
\end{enumerate}

\subsubsection{Layer structure used for the order-of-magnitude cascade}
The representative stratification adopted for the quantitative closed estimate is:
\begin{table}[ht]
\centering
\caption{Representative order-of-magnitude stratification used for the cascade narrative (not a substitute for measured profiles). ``Constr.'' abbreviates combined constraint strength from gravity and magnetic field in the text.}
\label{tab:iono-strata}
\renewcommand{\arraystretch}{1.2}
\scriptsize
\begin{tabular}{L{1.0cm} L{2.3cm} L{2.0cm} L{1.6cm} L{3.1cm} L{2.3cm}}
\hline
Region & Altitude range & $n$ ($\mathrm{m^{-3}}$) & $\sigma$ ($\mathrm{S/m}$) & Dominant carriers & \shortstack{\scriptsize Constr.} \\
\hline
D  & $60$--$90~\mathrm{km}$ & $10^8$--$10^9$ & $10^{-8}$--$10^{-7}$ & neutrals with small ion fraction & weak \\
E  & $90$--$150~\mathrm{km}$ & $10^9$--$10^{10}$ & $10^{-7}$--$10^{-6}$ & electrons/ions increasing & medium \\
F1 & $150$--$250~\mathrm{km}$ & $10^{10}$--$10^{11}$ & $10^{-6}$--$10^{-5}$ & electron/ion dominated & strong \\
F2 & $250$--$1000+$~km & $10^{11}$--$10^{12}$ & $10^{-5}$--$10^{-4}$ & fully ionized & very strong \\
\hline
\end{tabular}
\normalsize
\end{table}
These are standard order-of-magnitude expectations in ionospheric physics \cite{Kelley2009EarthIonosphere,SchunkNagy2009IonospherePhysics}.

\subsubsection{Dual asymmetry unification}
Gravity filtering selects high-energy carriers and seeds an initial weak charge separation (seed electric field). The magnetic constraint enforces directional motion and converts this separation into a directed drive current. The stratified ionosphere then amplifies both the field and current through multi-layer coupling. Electron--ion mass disparity further stabilizes the charge separation. Finally, constraint-induced statistical asymmetry supports a self-sustaining regime without continuous external power in the ideal limit.

Within the thermal working hypothesis, the same ingredients appear as follows: gravity supplies vertical speed--altitude sorting; multi-species collisions supply stepwise upward redistribution of kinetic energy among ions; the magnetic field limits cross-field dispersal and favors persistence of layer-like structure in which successive collisional steps may occur. The term ``dual asymmetry'' thus encompasses both rectified electrodynamics and a biased kinetic pathway capable, under suitable composition and collisionality, of sustaining hot, light-ion-enriched topside populations.

\section{Quantitative electrostatic estimate (seed field to macroscopic field)}
\subsection{Electric-field amplitudes: closure scale versus typical ionospheric reports}
\label{sec:efield-scales}
Large-scale electric fields in the terrestrial ionosphere are routinely discussed at $\mathrm{mV\,m^{-1}}$ to $\mathrm{V\,km^{-1}}$ scales in textbooks and monographs \cite{Kelley2009EarthIonosphere,SchunkNagy2009IonospherePhysics}. The algebraic closure in this section therefore uses an \emph{illustrative} target amplitude $E_{\mathrm{total}}\sim 10^2\ \mathrm{V\,m^{-1}}$ as a dimensionally consistent way to exhibit a $10^6$ multiplicative stacking of a seed field $E_0\sim 10^{-4}\ \mathrm{V\,m^{-1}}$. The same recursion relations apply if all electric fields are scaled uniformly by a common factor $\alpha$ (e.g., choosing $\alpha\sim 10^{-3}$ would map the narrative onto $\sim 10^{-1}\ \mathrm{V\,m^{-1}}$ while preserving the dimensionless factors $\mathcal{A}_i$). The physical identification of $E$ with a specific observable (volume-averaged DC field, effective field along a flux tube, or a localized structure) must be fixed when comparing to data.

\subsection{Basic parameters (reference scales)}
Representative terrestrial ionospheric scales and fundamental constants are taken as follows:
\begin{itemize}
  \item reference closure temperature: $T_{\mathrm{ref}}\approx \SI{300}{K}$;
  \item thermal voltage (closure scale): $V_T = kT_{\mathrm{ref}}/e \approx \SI{0.026}{V}$;
  \item electron number density in F2: $n\approx 10^{11}~\mathrm{m^{-3}}$;
  \item gravity scale: $g\approx \SI{9.8}{m/s^2}$;
  \item elementary charge $e$ and Boltzmann constant $k$ taken from CODATA \cite{Tiesinga2021CODATA};
  \item vacuum permittivity $\varepsilon_0$ taken from CODATA \cite{Tiesinga2021CODATA};
  \item electron and representative ion masses (mass ratio used as a scale) from CODATA \cite{Tiesinga2021CODATA};
  \item representative region conductivities: $\sigma_D\sim 10^{-8}$, $\sigma_E\sim 10^{-7}$, $\sigma_{F1}\sim 10^{-6}$, $\sigma_{F2}\sim 10^{-5}~\mathrm{S/m}$.
\end{itemize}
The numerical closure also introduces effective, electrostatic parameters $\eta$, $\xi$, and $\zeta$ for the mapping between micro-level constraint bias and macroscopic fields. Their values are constrained by consistency with the target field magnitude and Gauss-law order estimates. Unless explicitly stated otherwise, numbers in this section are order-of-magnitude illustrative values. Here $T_{\mathrm{ref}}$ is a closure reference temperature used in the electrostatic seed estimate; it is not identical to the elevated electron temperature used in the suprathermal-tail discussion.

\subsection{Seed field from Debye screening and constraint-to-electrostatics mapping}
In an isothermal plasma approximation, the Debye length is
\begin{equation}
\lambda_{\mathrm{D}}=\sqrt{\frac{\varepsilon_0 kT_{\mathrm{ref}}}{n e^2}}\ ,
\end{equation}
with the Debye-length form standard in plasma physics \cite{ChenPlasmaPhysics1984} and the constants taken from CODATA \cite{Tiesinga2021CODATA}. Using $T_{\mathrm{ref}}\approx \SI{300}{K}$ and $n\approx 10^{11}~\mathrm{m^{-3}}$ gives $\lambda_D\approx 3.78\times 10^{-3}~\mathrm{m}$.

We define an effective separation length
\begin{equation}
\ell_0=\xi \lambda_D,
\end{equation}
with $\xi\approx 8$.
We map the constraint-induced statistical bias to an electrostatic potential difference
\begin{equation}
\Delta\phi_0=\eta\,V_T,\qquad V_T=\frac{kT_{\mathrm{ref}}}{e},
\end{equation}
with $\eta=10^{-4}$.
Then the seed electric field is
\begin{equation}
E_0=\frac{\Delta\phi_0}{\ell_0}=\frac{\eta V_T}{\xi\lambda_D}\approx 8.5\times10^{-5}~\mathrm{V/m}.
\end{equation}

\subsubsection{Uncertainty propagation at closure level}
From $E_0\propto \eta\,T_{\mathrm{ref}}^{1/2}\,\xi^{-1}\,n^{1/2}$ (via $\lambda_D\propto\sqrt{T_{\mathrm{ref}}/n}$), first-order relative uncertainty gives
\begin{equation}
\frac{\delta E_0}{E_0}\approx
\frac{\delta\eta}{\eta}
+\frac{1}{2}\frac{\delta T_{\mathrm{ref}}}{T_{\mathrm{ref}}}
+\frac{1}{2}\frac{\delta n}{n}
-\frac{\delta\xi}{\xi}.
\end{equation}
This relation identifies which measurements or priors most strongly control the inferred seed-field uncertainty.

\subsection{Seed current density from an effective Ohm closure}
In the magnetic rectification stage, the directed current density is estimated via an effective scalar Ohm relation (a pedagogical reduction of tensor conductivity and Hall/Pedersen structure at finite collision frequency \cite{Kelley2009EarthIonosphere,SchunkNagy2009IonospherePhysics}):
\begin{equation}
J_0\approx \sigma_{\mathrm{eff}}E_0,
\end{equation}
where we identify $\sigma_{\mathrm{eff}}$ with the F2-scale conductivity $\sigma_{F2}$ adopted in the parameter list. With $\sigma_{F2}\approx 10^{-5}~\mathrm{S/m}$ one obtains
\begin{equation}
J_0\approx 8.5\times10^{-10}~\mathrm{A/m^2}.
\end{equation}

\subsection{Cascade amplification and macroscopic field}
Cascade amplification is modeled as a multiplicative layer-to-layer enhancement:
\begin{equation}
E_i=\mathcal{A}_i\,E_{i-1},
\end{equation}
with $i\in\{\mathrm{E},\mathrm{F1},\mathrm{F2}\}$.
To match the target macroscopic field scale, we adopt effective multipliers
\begin{equation}
\mathcal{A}_{\mathrm{E}}\approx 10,\qquad
\mathcal{A}_{\mathrm{F1}}\approx 1000,\qquad
\mathcal{A}_{\mathrm{F2}}\approx 100,
\end{equation}
so that the effective total amplification is
\begin{equation}
\mathcal{A}_{\mathrm{tot}}\approx 10\times 1000\times 100 = 10^6.
\end{equation}
Here $\mathcal{A}_{\mathrm{E}},\mathcal{A}_{\mathrm{F1}},\mathcal{A}_{\mathrm{F2}}$ are order-of-magnitude effective gains consistent with a layer-to-layer recursion picture; residual non-ideal loss/projection effects are compressed into $\zeta$. With $\zeta\approx 0.94$, the resulting macroscopic field becomes
\begin{equation}
E_{\mathrm{total}}\approx \zeta\,E_{\mathrm{F2}}\approx \SI{80}{V/m},
\end{equation}
which matches the illustrative closure scale adopted above; rescaling $\eta,\xi,\zeta$ (or an overall factor $\alpha$) maps the same recursion onto smaller amplitudes consistent with typical reported ionospheric fields \cite{Kelley2009EarthIonosphere,SchunkNagy2009IonospherePhysics}.

\subsection{Gauss-law self-consistency check}
A planar order-of-magnitude estimate is used to verify that the implied charge imbalance is not orders of magnitude larger than the nominal carrier density:
\begin{equation}
\Sigma_q\sim c\,\varepsilon_0 E_{\mathrm{total}},\qquad
\Delta n \sim \frac{\Sigma_q}{e\,\ell_0}\approx \frac{c\,\varepsilon_0 E_{\mathrm{total}}}{e\,\ell_0},
\end{equation}
where $c=O(1)$ accounts for geometric factors omitted in this order-of-magnitude planar estimate. Using $c\approx 1$, $\ell_0=\xi\lambda_D\approx 8\times 3.78\times 10^{-3}~\mathrm{m}\approx 3.0\times 10^{-2}~\mathrm{m}$ and $E_{\mathrm{total}}\approx \SI{80}{V/m}$ yields $\Delta n\approx 1.5\times 10^{11}~\mathrm{m^{-3}}$, comparable to the assumed $n\approx 10^{11}~\mathrm{m^{-3}}$.
This supports overall electrostatic self-consistency at the order-of-magnitude level.

\section{Discussion}
\subsection{Closed electrostatic parameterization and limitations}
The electrodynamic component is constructed as a self-contained closure. The effective parameters $\eta,\xi,\zeta$ map constraint-induced microscale bias to macroscopic electrostatic fields and compress non-ideal effects associated with layer-to-layer coupling.

These parameters should ultimately be obtained from constrained transport models (kinetic or fluid-kinetic) that specify collision operators, boundary conditions, and magnetic geometry explicitly. The present stage establishes \emph{existence and order-of-magnitude consistency} of the closure under minimal assumptions.

\paragraph{Thermodynamic consistency.}
Directed currents dissipate energy in resistive media; sustaining a steady state in a resistive plasma generally requires power balance with sources and/or boundary fluxes. The ``no external driving'' idealization is therefore a \emph{limiting statement concerning absent macroscopic thermal and chemical-potential gradients and absent continuous non-thermal injection} (Sec.~\ref{sec:intro}), not an assertion that the terrestrial ionosphere constitutes a closed equilibrium isolated from the Sun. In applications, residual photoionization, neutral winds, and magnetospheric coupling should be represented by explicit source terms when comparing to observations.

\paragraph{Quasineutrality and the Debye closure.}
Plasmas are quasineutral at scales $\gg \lambda_D$; the use of $\lambda_D$ here fixes the transverse scale of charge-separation fluctuations that can seed an electrostatic potential, not a macroscopic net charge density comparable to $n$ over the entire layer volume.

\subsection{Relation to traditional ionospheric mechanisms}
The formulation does not assert that external drivers are negligible in real ionospheres. It isolates an intrinsic, constraint-driven pathway that remains formally defined in the ideal limit of vanishing macroscopic temperature gradients and external non-thermal injection. Quantitative comparison with observations requires inclusion of residual external sources and characterization of their effect on the effective closure parameters.

\subsection{Correspondence with statistical physics: constraint reshapes occupation distributions}
According to \cite{peng2026entropydirectionmirrorstateparadox}, constraints can reshape the long-time occupation distribution such that directed transport arises without macroscopic thermal gradients. Gravitational selection and magnetic rectification supply a concrete planetary realization of this constraint-to-transport mapping. Cascade amplification in a stratified ionosphere connects microscopic bias to macroscopic fields.

\subsection{Testable trends}
The parameterized formulation yields the following qualitative dependences:
\begin{enumerate}
  \item $E_0=\eta V_T/(\xi\lambda_D)$, with $\lambda_D\propto \sqrt{T_{\mathrm{ref}}/n}$, yields explicit dependence of the seed field on $T_{\mathrm{ref}}$ and $n$;
  \item $E_{\mathrm{total}}$ depends on the effective layer multipliers, which in turn reflect conductivity stratification and coupling efficiency;
  \item electron--ion mass disparity influences the stability of charge separation and hence the persistence of directed currents;
  \item in a coarse-grained representation of statistically rare multi-collision sequences, the suprathermal light-ion fraction scales as $f_{\mathrm{hot}}\sim p^N$, linking collision statistics to upper-tail observables.
\end{enumerate}

\subsection{High effective temperatures in ionospheres and corona-like environments}
Conventional heating models emphasize external energy deposition and balance with radiative and collisional cooling, thereby addressing temperature structure under specified composition and loss rates. Within the present constraint-based framework, large effective temperatures may alternatively be interpreted as dynamic enrichment of kinetic energy maintained by \emph{speed-space filtering} and \emph{multi-species collisional redistribution}.

In a single-species plasma, collisions primarily redistribute momentum within one mass species and are less effective at sustaining a multi-stage channel for selective energization. When light and heavy ions coexist, elastic collisions can repeatedly transfer energy to light ions in favorable geometries: following gravitational filtering, a small subpopulation may traverse a sequence of encounters with heavy partners that successively increase kinetic energy relative to the bulk. The result is a suprathermal tail enriched in light ions (in particular protons and hydrogenic species).

\subsubsection{Light-ion energy gain in collisions with heavy ions: one-dimensional estimate}
For illustration, consider a one-dimensional elastic collision between a light ion ($m_\ell$, e.g., H$^+$) and a heavy ion ($m_h\gg m_\ell$) with pre-collision velocities $u_\ell,u_h$. The post-collision light-ion velocity is
\begin{equation}
v_\ell'=\frac{m_\ell-m_h}{m_\ell+m_h}u_\ell+\frac{2m_h}{m_\ell+m_h}u_h.
\end{equation}
For $m_h\gg m_\ell$, this reduces to
\begin{equation}
v_\ell'\approx -u_\ell+2u_h.
\end{equation}
Hence the light ion gains kinetic energy whenever $|v_\ell'|>|u_\ell|$, i.e.
\begin{equation}
|{-}u_\ell+2u_h|>|u_\ell|.
\end{equation}
In head-on geometry with $u_h<0<u_\ell$, the gain condition is naturally satisfied and the heavy ion acts as an effective moving reflector. This is a standard elastic-collision consequence of mass asymmetry, not an external energy source \cite{ChenPlasmaPhysics1984}.

In three dimensions, collision angles and velocity orientations are distributed, and only a fraction of encounters increase H$^+$ kinetic energy. A coarse-grained representation replaces the microscopic distribution by an effective single-step gain factor,
\begin{equation}
E_{k+1}=g\,E_k,\qquad g>1
\end{equation}
for favorable collisions, while unfavorable events return ions to the bulk.

\subsubsection{Multi-step energization: order-of-magnitude illustration}
\label{sec:multistep-ion}
Consider an H$^+$ ion with initial kinetic energy $E_0$ and suppose that $N$ successive favorable collisions are required to exceed an upper-layer threshold $E_{\mathrm{th}}$:
\begin{equation}
E_0\,g^N\ge E_{\mathrm{th}}
\quad\Rightarrow\quad
N\ge \frac{\ln(E_{\mathrm{th}}/E_0)}{\ln g}.
\end{equation}
Using illustrative values $E_0=1\,\mathrm{eV}$, $E_{\mathrm{th}}=30\,\mathrm{eV}$, and $g=1.5$, one gets
\begin{equation}
N\approx \frac{\ln 30}{\ln 1.5}\approx 8.4,
\end{equation}
so approximately nine favorable steps are required.

If each step is favorable with probability $p$, the probability of an uninterrupted favorable chain of length $N$ is approximately
\begin{equation}
P_N\sim p^N.
\end{equation}
For $p=0.3$ and $N=9$, $P_N\sim 2\times10^{-5}$. The smallness of $P_N$ indicates that only a minute subpopulation of H$^+$ ions attains the highest effective temperatures at upper altitudes, while the bulk remains near lower energies.

The corresponding one-particle equivalent temperature scale can be expressed as
\begin{equation}
T_{\mathrm{eq}}=\frac{2E}{3k}.
\end{equation}
For $E_{\mathrm{th}}=30\,\mathrm{eV}$, this yields $T_{\mathrm{eq}}\approx 2.3\times 10^5\,\mathrm{K}$, indicating that the schematic channel generates a high-energy \emph{tail} rather than a uniformly heated Maxwellian bulk.

If $N_{\ell}$ denotes the number of candidate light ions entering the cascade region over a characteristic residence time, the expected number of high-tail ions is
\begin{equation}
N_{\mathrm{hot}}\sim N_{\ell}P_N.
\end{equation}
For example, $N_{\ell}\sim 10^{10}$ with $P_N\sim 2\times10^{-5}$ yields $N_{\mathrm{hot}}\sim 2\times10^5$: a small suprathermal population that may nevertheless dominate certain high-energy diagnostics.

\paragraph{Limitations of the coarse-grained estimate.}
The relation $P_N\sim p^N$ presumes approximately independent effective gain events and a single representative gain factor $g$. Ionospheric Coulomb transport is in reality dominated by many small-angle deflections and pitch-angle scattering; $(g,p)$ should therefore be regarded as effective parameters subject to calibration by kinetic or Fokker--Planck models \cite{ChenPlasmaPhysics1984,Kelley2009EarthIonosphere,SchunkNagy2009IonospherePhysics}. In correlated collisional environments, the effective number of independent steps (or an effective $p_{\mathrm{eff}}$) is reduced; $P_N$ serves as an order-of-magnitude surrogate rather than as a closed-form probability.

Electrons may possess thermal speeds large compared with nominal gravitational escape scales yet remain confined: the positive-ion population establishes an electrostatic potential that limits net electron loss, and magnetic geometry contributes to circulation and confinement consistent with replenishment.

Composition-dependent trends across planetary and coronal environments are therefore stated at the level of \emph{qualitative predictions}: a substantial heavy-ion component may favor multi-species collision pathways that sustain energetic light-ion tails, whereas an ionosphere composed almost entirely of hydrogen and helium may provide fewer heavy partners for repeated mass-asymmetric transfer. Quantitative application requires measured composition, realistic collision operators, and radiative cooling; inter-object comparisons in the present text are not asserted as established empirical correlations.

\subsection{Comparative implications: Venus, Jupiter, hydrogenic topside, and magnetic confinement}
The working hypothesis is constructed to subsume several stylized observational regularities within a single framework, without independent \emph{ad hoc} explanations for each environment. The unifying structure comprises gravitational speed--altitude sorting, the ion mass spectrum as it controls multi-step collisional redistribution, and magnetic stabilization of stratified plasma where the field is strong.

\paragraph{Venus.}
Under heavy-ion dominance and limited light-ion abundance, the multi-step collisional sequence that energizes a minority population to very high speeds possesses both fewer effective stages and fewer candidate light ions: collisions remain frequent, but the mass-asymmetric channel lacks the light-ion reservoir required to build a thick suprathermal tail at altitude. In the absence of a strong global magnetic field, large-scale plasma structure is more susceptible to scattering and cross-field dilution. Within the present hypothesis, the comparatively weak Venus ionosphere is attributed in part to the conjunction of composition-limited upward kinetic redistribution and reduced magnetically supported stratification; solar EUV heating and ionization are not dismissed and remain necessary ingredients in any complete model.

\paragraph{Jupiter.}
When the ionized population consists predominantly of hydrogen and helium, the number of distinct heavy-ion partners available as momentum reservoirs in elastic encounters is restricted. Trace heavy ions may be present, but if their abundance or collision rate with the bulk is small, the multi-step mechanism is less efficient at transferring energy repeatedly to the dominant light species. The framework therefore permits comparatively moderate topside thermal enhancement relative to environments with abundant oxygen, molecular ions, or other heavy species at similar external EUV forcing, pending detailed numerical assessment.

\paragraph{Hydrogenic dominance at the topside.}
Among common ionospheric constituents, hydrogen (protons) has the smallest mass and, for specified momentum transfer in favorable collisions, the largest velocity increment. In combination with gravitational filtering, the hypothesis implies that the highest-altitude, hottest \emph{ionic} component should frequently be hydrogenic when hydrogen is available: it is the species most readily elevated through the collisional sequence and, at given kinetic energy, exhibits the largest gravitational scale height. The well-documented terrestrial transition to H$^+$-rich topside plasma is thus interpreted here as a consequence of the same structure rather than as an independent empirical regularity.

\paragraph{Vertical mass distribution and maximum effective temperature.}
The peak effective temperature (or suprathermal-tail temperature) attainable at the column top is not universal: it should scale with the efficiency of upward kinetic-energy transfer permitted by the local species mix. A qualitative control parameter is the effective availability of heavy-ion collision partners for light ions over altitudes where the collision frequency remains sufficient to sustain many successive steps. Variation of relative abundances between layers alters the effective step count and the probability $P_N$ of long favorable sequences (Sec.~\ref{sec:multistep-ion}). The vertical distribution of ion masses therefore belongs to the prediction space of the hypothesis together with EUV heating and radiative cooling.

\paragraph{Role of the magnetic field.}
Even when gravity and collisions supply energy to a light-ion tail, the large-scale \emph{intensity} of an ionosphere or corona-like envelope depends additionally on whether the plasma retains organized structure over times long enough for stratified amplification or repeated collisional steps. Magnetic geometry inhibits rapid cross-field loss of electrons, channels currents, and may suppress turbulent mixing that would erode sharp layers. The field does not substitute for external heating but may lengthen the effective coupling distance over which the hypothesized cascade operates.

\paragraph{Synthesis.}
The working hypothesis advanced here supplies a single structural basis for (i) elevated temperatures in ionospheres and coronae above cooler dense photospheres or surfaces, (ii) vertical species sorting with light ions preferentially found at large altitude, (iii) systematic inter-object differences governed by ion mass spectra and magnetic environment, and (iv) electrodynamic signatures tied to the same dual constraints. Quantitative ordering of environments (Venus versus Jupiter, Earth versus the solar corona) requires kinetic simulations with measured composition; the preceding discussion establishes qualitative causal relations only.

\subsection{Outlook on quantitative validation}
The preceding analysis is qualitative. A necessary extension is the incorporation of radiative cooling and realistic collision cross sections into constrained transport models, together with derivation of bounds on $\eta,\xi,\zeta$ from first principles or numerical kinetic simulation rather than phenomenological assignment alone.

\section{Conclusion and outlook}
This paper has presented a gravity--magnetic dual-asymmetry formulation for unbiased rectification and cascade amplification in planetary ionospheres. The central thermal claim is a \emph{working hypothesis}: ionospheres and coronae may exhibit effective temperatures greatly exceeding those of the underlying surface or photosphere as a consequence of gravitational sorting of faster particles to greater altitude, multi-step elastic collisions between light and heavy ions that redistribute kinetic energy upward, and magnetic stabilization of stratified structure where the field is strong. In parallel, within the statistical-mechanical principle that asymmetric constraints reshape steady occupation distributions, gravity supplies speed-space filtering and the magnetic field rectifies seed charge separation into directed current; the stratified ionosphere closes the electrostatic recursion at order of magnitude. An illustrative field amplitude $E\sim 10^2\ \mathrm{V\,m^{-1}}$ was exhibited, with the caveat that typical reported large-scale ionospheric fields are much smaller unless the closure is rescaled uniformly (Sec.~\ref{sec:efield-scales}).

The same working hypothesis was applied qualitatively to relate terrestrial topside H$^+$ enrichment, conditions approximating Venus (heavy-ion-rich plasma, weak intrinsic field), and Jupiter-like hydrogen--helium ionospheres (limited heavy-ion diversity) as distinct realizations of one structural mechanism controlled by composition and magnetism.

Recommended extensions include (i) derivation of $\eta,\xi,\zeta$ within constrained kinetic or fluid-kinetic transport models, (ii) replacement of effective recursion factors by expressions tied explicitly to conservation laws and boundary conditions, (iii) systematic comparison with observational data sets including scaling and uncertainty analysis across geophysical conditions, and (iv) quantitative assessment of the thermal unification using multi-fluid or kinetic models with species-resolved collision operators and radiative losses.

\bibliographystyle{plain}
\bibliography{references}

\end{document}

